
\documentclass[a4paper,12pt]{article}
\usepackage[italian]{babel}
\usepackage{geometry}
\usepackage{listings}
\geometry{margin=2.5cm}

\title{Verifica SQL - Sistema Universitario}
\date{}

\begin{document}
\maketitle

Si consideri il seguente schema logico relativo ad un sistema informativo universitario.

STUDENTE(\underline{idStudente}, nome, cognome, citta, annoIscrizione)

DOCENTE(\underline{idDocente}, nome, cognome, dipartimento)

CORSO(\underline{idCorso}, nomeCorso, crediti, idDocente)

AULA(\underline{idAula}, edificio, capienza)

ISCRIZIONE(\underline{idIscrizione}, idStudente, idCorso, voto, dataEsame)

\section*{Query SQL}

Scrivere le seguenti interrogazioni SQL.



\begin{enumerate}
	
	\item Elencare nome e cognome degli studenti ordinati per cognome e, a parità di cognome, per nome.
	
	\item Elencare nome del corso, crediti e cognome del docente che lo tiene, ordinando i risultati per crediti decrescenti.
	
	\item Elencare nome e cognome dei docenti appartenenti al dipartimento di Informatica che tengono almeno un corso.
	
	\item Elencare nome, cognome e annoIscrizione degli studenti residenti nella città di Roma ordinati per anno di iscrizione.
	
	\item Elencare nomeCorso e crediti dei corsi con più di 6 crediti indicando anche il cognome del docente responsabile.
	
	\item Elencare nome, cognome dello studente e nome del corso a cui è iscritto.
	
	\item Elencare nome dello studente, nome del corso e voto ottenuto per tutti gli esami sostenuti.
	
	\item Contare quanti studenti risultano iscritti a ciascun corso mostrando nomeCorso e numero di studenti.
	
	\item Calcolare per ogni corso il voto medio degli esami sostenuti mostrando nomeCorso e voto medio.
	
	\item Elencare nome e cognome degli studenti che non hanno mai sostenuto alcun esame.
	
	\item Elencare i corsi che non hanno ancora studenti iscritti.
	
	\item Trovare i docenti che tengono più di un corso mostrando cognome del docente e numero di corsi.
	
	\item Trovare gli studenti che hanno sostenuto almeno due esami mostrando nome, cognome e numero di esami sostenuti.
	
	\item Elencare i corsi il cui voto medio degli esami è maggiore di 25 mostrando nomeCorso e voto medio.
	
	\item Trovare lo studente che ha sostenuto il maggior numero di esami mostrando idStudente e numero di esami.
	
	\item Elencare i docenti i cui corsi hanno almeno tre studenti iscritti mostrando cognome del docente e numero di studenti complessivi.
	
	\item Elencare i corsi che non hanno mai avuto esami registrati.
	
	\item Trovare il corso con il voto medio più alto mostrando nomeCorso e voto medio.
	
	\item Elencare nome e cognome degli studenti che hanno ottenuto almeno un voto pari a 30.
	
	\item Elencare nome e cognome degli studenti che hanno sostenuto esami in almeno due corsi diversi.
	
\end{enumerate}


\end{document}
